\chapter{מבוא לדוגמנות שפה גדולה}

\section{רקע ודוגמנות שפה}

דוגמנות שפה היא תחום בעל חשיבות ליסודית בעיבוד שפה טבעית. לאורך שנים רבות,
מודלים סטטיסטיים שונים שימשו לחיזוי האסימון הבא בתוך רצף נתון של אסימונים.
עם התפתחות הלמידה העמוקה, חלה מהפכה בדרך בה אנו מעצבים מערכות אלה.

\section{רשתות עצביות וללמידה עמוקה}

רשתות עצביות מלאכותיות הן מודלים חישוביים המושראים מן המוח הביולוגי.
כל רשת מורכבת מ\textenglish{neurons} וממשקלים המחברים ביניהם.
ההדרכה של רשתות אלה נעשית באמצעות אלגוריתם בשם \textenglish{backpropagation},
אשר מחשב את הנגזרות הדרושות לעדכון המשקלים.

\section{אתגרים בדוגמנות של רצפים ארוכים}

אחד האתגרים העיקריים בעבודה עם רצפים טקסט ארוכים הוא זכירת תלויות ארוכות טווח.
רשתות \textenglish{RNN} סבלו מבעיית ההיעלמות או הפיצוץ של הגרדיאנט.
רשתות \textenglish{LSTM} שיפרו זאת מעט, אך עדיין לא היו מספיקות לטקסטים שלמים או מסמכים גדולים.

\section{המהלך לכיוון Attention}

מנגנון תשומת הלב החל להשתמש בתחום הראייה הממוחשבת, כדי לתרגם תמונות.
בעבודות מאוחרות יותר, חוקרים הבינו שניתן להחיל תשומת הלב גם על בעיות בשפה טבעית.
הרעיון הבסיסי הוא שהמודל צריך ללמוד איזה חלקים מהקלט רלוונטיים לכל פלט.

\section{המולד של \textenglish{Transformers}}

ב־2017 פורסמה המאמר \textenglish{Attention Is All You Need},
אשר הציג את ארכיטקטורת \textenglish{Transformer}.
מודל זה התבסס כולו על מנגנון תשומת הלב, ללא שימוש בשום \textenglish{RNN} או \textenglish{CNN}.
זה אפשר למודלים בעלי פרמטרים רבים בהרבה לתרגל בצורה יעילה ובמקביל.

\section{השפעתם של דוגמנות שפה גדולה}

בשנים האחרונות, דוגמנות שפה גדולה כגון \textenglish{GPT} ו־\textenglish{BERT}
הפכו לכלים בסיסיים ברבים מתחומי הבינה המלאכותית.
מודלים אלה מאומנים על כמויות ענקיות של טקסט, ויכולים לטפל במגוון רחב של משימות.
הם השפיעו על כל ההנדסה של \textenglish{NLP} וגם על תחומים סמוכים.
